% Blueprint: Minkowski-Caratheodory theorem
% Created by Numina

\begin{document}

\section{Minkowski--Carathéodory Theorem}

Throughout, $E$ is a finite-dimensional real normed space, $s \subseteq E$ is a
compact convex set, and $x \in s$. We write $\operatorname{ext}(s)$ for the set
of extreme points of $s$ (Mathlib: \texttt{Set.extremePoints}). The goal is the
\emph{Minkowski--Carathéodory theorem}: every point of $s$ is a convex
combination of at most $\operatorname{finrank}_{\mathbb R} E + 1$ extreme points
of $s$.

The proof combines two ingredients: a cardinality-bounded form of
Carathéodory's theorem, and \emph{Minkowski's theorem} (the finite-dimensional
Krein--Milman theorem), which states that $s$ equals the convex hull --- not
merely its closure --- of its extreme points.

\subsection{Carathéodory with a dimension bound}

\begin{lemma}[Carathéodory representation lands in the convex hull of a finite image]
    \label{lem:caratheodory-image-mem}
    \lean{LeanEval.ConvexGeometry.caratheodory_image_mem}
    \leanfile{LeanEval/ConvexGeometry/MinkowskiCaratheodory.lean}
    \leanok
    Let $A \subseteq E$ and let $x \in \operatorname{convexHull}_{\mathbb R} A$.
    Then there are a fintype $\iota$, an affinely independent family
    $z : \iota \to E$ with $\operatorname{range} z \subseteq A$, and strictly
    positive weights $w : \iota \to \mathbb R$ with $\sum_i w_i = 1$ and
    $\sum_i w_i\, z_i = x$, such that, writing $t = z\,''\,\iota$ for the image
    finset, one has $t \subseteq A$ and
    $x \in \operatorname{convexHull}_{\mathbb R} t$.
\end{lemma}
\begin{proof}
    By the explicit Carathéodory representation
    (\texttt{eq\_pos\_convex\_span\_of\_mem\_convexHull}) we obtain $\iota$, $z$,
    and $w$ as stated, with $\operatorname{range} z \subseteq A$. Let $t$ be the
    image of $z$; then $t \subseteq A$, and $x = \sum_i w_i\, z_i$ is the center of
    mass of the points $z_i$ with strictly positive weights summing to $1$, hence
    a convex combination of the points of $t$, so
    $x \in \operatorname{convexHull}_{\mathbb R} t$
    (\texttt{Finset.centerMass\_id\_mem\_convexHull}).
\end{proof}

\begin{lemma}[Carathéodory cardinality estimate]
    \label{lem:caratheodory-card-bound}
    \lean{LeanEval.ConvexGeometry.caratheodory_card_bound}
    \leanfile{LeanEval/ConvexGeometry/MinkowskiCaratheodory.lean}
    \uses{lem:caratheodory-image-mem}
    \leanok
    With $\iota$, $z$, and $t = z\,''\,\iota$ as in
    Lemma~\ref{lem:caratheodory-image-mem}, one has
    \[
        \#t \;\le\; \#\iota \;\le\; \operatorname{finrank}_{\mathbb R} E + 1 .
    \]
\end{lemma}
\begin{proof}
    \uses{lem:caratheodory-image-mem}
    Since $t = z\,''\,\iota$ is the image of $\iota$ under $z$, we have
    $\#t \le \#\iota$ (\texttt{Finset.card\_image\_le}). As $z$ is affinely
    independent,
    $\#\iota \le \operatorname{finrank}_{\mathbb R}\big(\operatorname{vectorSpan}_{\mathbb R}(\operatorname{range} z)\big) + 1$
    (\texttt{AffineIndependent.card\_le\_finrank\_succ}). Finally
    $\operatorname{vectorSpan}_{\mathbb R}(\operatorname{range} z)$ is a submodule of
    $E$, so its \texttt{finrank} is at most $\operatorname{finrank}_{\mathbb R} E$
    (\texttt{Submodule.finrank\_le}); combining the three inequalities gives
    $\#t \le \#\iota \le \operatorname{finrank}_{\mathbb R} E + 1$.
\end{proof}

\begin{lemma}[Carathéodory cardinality bound]
    \label{lem:caratheodory-card}
    \lean{LeanEval.ConvexGeometry.caratheodory_card}
    \leanfile{LeanEval/ConvexGeometry/MinkowskiCaratheodory.lean}
    \uses{lem:caratheodory-image-mem, lem:caratheodory-card-bound}
    \leanok
    Let $A \subseteq E$ and let $x \in \operatorname{convexHull}_{\mathbb R} A$.
    Then there is a finite set $t \subseteq A$ with
    $\#t \le \operatorname{finrank}_{\mathbb R} E + 1$ and
    $x \in \operatorname{convexHull}_{\mathbb R} t$.
\end{lemma}
\begin{proof}
    \uses{lem:caratheodory-image-mem, lem:caratheodory-card-bound}
    Take $t = z\,''\,\iota$ from Lemma~\ref{lem:caratheodory-image-mem}; then
    $t \subseteq A$ and $x \in \operatorname{convexHull}_{\mathbb R} t$. The
    cardinality bound $\#t \le \operatorname{finrank}_{\mathbb R} E + 1$ is
    Lemma~\ref{lem:caratheodory-card-bound}.
\end{proof}

\subsection{Building blocks for Minkowski's theorem}

For a continuous linear functional $\ell : E \to \mathbb R$ we single out the
face of $s$ on which $\ell$ attains its maximum.

\begin{definition}[Maximal face of a linear functional]
    \label{def:exposed-face}
    \lean{LeanEval.ConvexGeometry.exposedFace}
    \leanfile{LeanEval/ConvexGeometry/MinkowskiCaratheodory.lean}
    \leanok
    For a continuous linear functional $\ell : E \to \mathbb R$ set
    \[
        F_\ell(s) \;=\; \{\, y \in s : \forall z \in s,\ \ell z \le \ell y \,\}.
    \]
    Equivalently, $F_\ell(s)$ is the exposed subset of $s$ determined by $\ell$
    (Mathlib: \texttt{IsExposed}).
\end{definition}

\begin{lemma}[The maximal face is an extreme subset]
    \label{lem:exposed-face}
    \lean{LeanEval.ConvexGeometry.exposedFace_isExtreme}
    \leanfile{LeanEval/ConvexGeometry/MinkowskiCaratheodory.lean}
    \uses{def:exposed-face}
    \leanok
    For every continuous linear functional $\ell$, the set $F_\ell(s)$ is compact
    and convex, and its extreme points are extreme points of $s$:
    $\operatorname{ext}(F_\ell(s)) \subseteq \operatorname{ext}(s)$.
\end{lemma}
\begin{proof}
    $F_\ell(s)$ is an exposed subset of $s$ (it is cut out by $\ell$), hence by
    \texttt{IsExposed.isExtreme} it is an extreme subset of $s$; therefore
    \texttt{IsExtreme.extremePoints\_subset\_extremePoints} gives
    $\operatorname{ext}(F_\ell(s)) \subseteq \operatorname{ext}(s)$. Compactness
    and convexity are \texttt{IsExposed.isCompact} (using compactness of $s$) and
    \texttt{IsExposed.convex}.
\end{proof}

\begin{lemma}[The maximal face lies in the kernel of $\ell$]
    \label{lem:face-vectorspan-le-ker}
    \lean{LeanEval.ConvexGeometry.exposedFace_vectorSpan_le_ker}
    \leanfile{LeanEval/ConvexGeometry/MinkowskiCaratheodory.lean}
    \uses{def:exposed-face}
    \leanok
    Every difference of two points of $F_\ell(s)$ is annihilated by $\ell$, hence
    \[
        \operatorname{vectorSpan}_{\mathbb R} F_\ell(s) \;\le\; \ker \ell .
    \]
\end{lemma}
\begin{proof}
    \uses{def:exposed-face}
    Any $y, y' \in F_\ell(s)$ both attain the maximal value of $\ell$ on $s$, so
    $\ell(y - y') = \ell y - \ell y' = 0$, i.e.\ $y - y' \in \ker \ell$. Thus every
    generator $y - y'$ of $\operatorname{vectorSpan}_{\mathbb R} F_\ell(s)$ lies in
    the submodule $\ker \ell$, so the whole span is contained in it
    (\texttt{vectorSpan\_le}).
\end{proof}

\begin{lemma}[The maximal face has smaller affine dimension]
    \label{lem:exposed-face-lower-dim}
    \lean{LeanEval.ConvexGeometry.exposedFace_finrank_lt}
    \leanfile{LeanEval/ConvexGeometry/MinkowskiCaratheodory.lean}
    \uses{def:exposed-face, lem:face-vectorspan-le-ker}
    \leanok
    Suppose $\ell$ is non-constant on $s$, i.e.\ there exist $a, b \in s$ with
    $\ell a \ne \ell b$. Then
    \[
        \operatorname{finrank}_{\mathbb R}\big(\operatorname{vectorSpan}_{\mathbb R} F_\ell(s)\big)
        \;<\;
        \operatorname{finrank}_{\mathbb R}\big(\operatorname{vectorSpan}_{\mathbb R} s\big).
    \]
\end{lemma}
\begin{proof}
    \uses{lem:face-vectorspan-le-ker}
    Write $U = \operatorname{vectorSpan}_{\mathbb R} F_\ell(s)$ and
    $V = \operatorname{vectorSpan}_{\mathbb R} s$. By
    Lemma~\ref{lem:face-vectorspan-le-ker}, $U \le \ker \ell$, and $U \le V$ since
    $F_\ell(s) \subseteq s$ (\texttt{vectorSpan\_mono}); hence
    $U \le V \sqcap \ker \ell$. As $a, b \in s$, the vector $a - b$ lies in $V$
    (\texttt{vsub\_mem\_vectorSpan}), whereas $\ell(a - b) = \ell a - \ell b \ne 0$
    shows $a - b \notin \ker \ell$. Therefore $V \sqcap \ker \ell$ is a
    \emph{proper} submodule of $V$, i.e.\ $V \sqcap \ker \ell < V$. Combining
    \[
        \operatorname{finrank}_{\mathbb R} U
        \;\le\; \operatorname{finrank}_{\mathbb R}(V \sqcap \ker \ell)
        \qquad (\texttt{Submodule.finrank\_mono})
    \]
    with
    $\operatorname{finrank}_{\mathbb R}(V \sqcap \ker \ell) < \operatorname{finrank}_{\mathbb R} V$
    (\texttt{Submodule.finrank\_lt\_finrank\_of\_lt}) yields the claim.
\end{proof}

To analyse relative boundary points we reduce to the \emph{direction subspace}.
Fix a base point $p \in s$ and let $W = \operatorname{vectorSpan}_{\mathbb R} s$,
which (as a submodule of the finite-dimensional space $E$) is itself a
finite-dimensional real normed space. Every difference of points of $s$ lies in
$W$, so $s \subseteq p + W$, and the map
$\iota_p : W \hookrightarrow E,\ w \mapsto w + p$ is an affine isometry. Write
\[
    s' \;=\; \{\, w \in W : w + p \in s \,\} \;=\; \iota_p^{-1}(s) \;\subseteq\; W
\]
for the translated copy of $s$ inside $W$; then $\iota_p\,''\,s' = s$ and
$0 \in s'$ (as $p \in s$).

\begin{lemma}[The translated set is convex]
    \label{lem:translated-set-convex}
    \lean{LeanEval.ConvexGeometry.translatedSet_convex}
    \leanfile{LeanEval/ConvexGeometry/MinkowskiCaratheodory.lean}
    \leanok
    The set $s' = \{\, w \in W : w + p \in s \,\}$ is convex.
\end{lemma}
\begin{proof}
    $s'$ is the preimage of the convex set $s$ under the affine map
    $w \mapsto w + p$, hence convex (\texttt{Convex.affine\_preimage}).
\end{proof}

We decompose the identity $\operatorname{intrinsicInterior}_{\mathbb R} s =
\iota_p\,''\,\operatorname{interior} s'$ into three independent pieces: the
translation recovers $s$, the translated set spans all of $W$, and (a general
fact) the intrinsic interior of a set whose affine span is the whole space is its
ordinary interior.

\begin{lemma}[The translation recovers $s$]
    \label{lem:iie-image}
    \lean{LeanEval.ConvexGeometry.image_translatedSet_eq}
    \leanfile{LeanEval/ConvexGeometry/MinkowskiCaratheodory.lean}
    \leanok
    For $p \in s$ and $s' = \{\, w \in W : w + p \in s \,\}$ one has
    $\iota_p\,''\,s' = s$, i.e.
    \[
        \{\, w + p : w \in W,\ w + p \in s \,\} \;=\; s .
    \]
\end{lemma}
\begin{proof}
    The inclusion $\iota_p\,''\,s' \subseteq s$ is immediate from the definition of
    $s'$. Conversely, for $z \in s$ the difference $z - p$ lies in
    $W = \operatorname{vectorSpan}_{\mathbb R} s$
    (\texttt{vsub\_mem\_vectorSpan}), and $\iota_p(z - p) = (z - p) + p = z \in s$,
    so $z \in \iota_p\,''\,s'$.
\end{proof}

\begin{lemma}[The translated set spans its subspace]
    \label{lem:iie-affinespan-top}
    \lean{LeanEval.ConvexGeometry.affineSpan_translatedSet_eq_top}
    \leanfile{LeanEval/ConvexGeometry/MinkowskiCaratheodory.lean}
    \uses{lem:iie-image}
    \leanok
    For $p \in s$ the translated set $s' = \{\, w \in W : w + p \in s \,\}$ has full
    affine span in $W$:
    \[
        \operatorname{affineSpan}_{\mathbb R} s' \;=\; \top .
    \]
\end{lemma}
\begin{proof}
    \uses{lem:iie-image}
    The set $s'$ is nonempty, as $0 \in s'$ (since $0 + p = p \in s$). The
    injective inclusion $\iota_p$ carries $\operatorname{vectorSpan}_{\mathbb R} s'$
    onto $\operatorname{vectorSpan}_{\mathbb R}(\iota_p\,''\,s')
    = \operatorname{vectorSpan}_{\mathbb R} s = W$ (using
    Lemma~\ref{lem:iie-image}), so $\operatorname{vectorSpan}_{\mathbb R} s'$ is all
    of $W$. A nonempty set whose vector span is the whole space has full affine
    span (\texttt{affineSpan\_eq\_top\_iff\_vectorSpan\_eq\_top\_of\_nonempty}).
\end{proof}

\begin{lemma}[Intrinsic interior equals interior under full affine span]
    \label{lem:intrinsicinterior-eq-interior}
    \lean{LeanEval.ConvexGeometry.intrinsicInterior_eq_interior_of_affineSpan_eq_top}
    \leanfile{LeanEval/ConvexGeometry/MinkowskiCaratheodory.lean}
    \leanok
    Let $F$ be a real normed space and $u \subseteq F$ with
    $\operatorname{affineSpan}_{\mathbb R} u = \top$. Then
    $\operatorname{intrinsicInterior}_{\mathbb R} u = \operatorname{interior} u$.
\end{lemma}
\begin{proof}
    By definition $\operatorname{intrinsicInterior}_{\mathbb R} u$ is the image,
    under the coercion $\operatorname{affineSpan}_{\mathbb R} u \hookrightarrow F$,
    of the topological interior of the preimage of $u$. When
    $\operatorname{affineSpan}_{\mathbb R} u = \top$ this coercion is the
    homeomorphism $\top \simeq F$ (\texttt{AffineSubspace.topHomeomorph} via
    \texttt{AffineSubspace.topEquiv}), under which preimage and interior commute,
    so the image is exactly $\operatorname{interior} u$.
\end{proof}

\begin{lemma}[Intrinsic interior via the direction subspace]
    \label{lem:intrinsic-interior-eq-image}
    \lean{LeanEval.ConvexGeometry.intrinsicInterior_eq_image}
    \leanfile{LeanEval/ConvexGeometry/MinkowskiCaratheodory.lean}
    \uses{lem:iie-image, lem:iie-affinespan-top, lem:intrinsicinterior-eq-interior}
    \leanok
    The affine span of $s'$ is all of $W$, so the intrinsic interior of $s'$ in
    $W$ coincides with its topological interior, and
    \[
        \operatorname{intrinsicInterior}_{\mathbb R} s
        \;=\; \iota_p\,''\,\operatorname{interior} s' .
    \]
\end{lemma}
\begin{proof}
    \uses{lem:iie-image, lem:iie-affinespan-top, lem:intrinsicinterior-eq-interior}
    The map $\iota_p : W \to E$, $w \mapsto w + p$, is an affine isometry (the
    isometric subspace inclusion followed by translation by $p$). By
    Lemma~\ref{lem:iie-image}, $\iota_p\,''\,s' = s$, so
    \texttt{AffineIsometry.image\_intrinsicInterior} gives
    \[
        \operatorname{intrinsicInterior}_{\mathbb R} s
        = \operatorname{intrinsicInterior}_{\mathbb R}(\iota_p\,''\,s')
        = \iota_p\,''\,\operatorname{intrinsicInterior}_{\mathbb R} s' .
    \]
    By Lemma~\ref{lem:iie-affinespan-top}, $\operatorname{affineSpan}_{\mathbb R} s'
    = \top$, so Lemma~\ref{lem:intrinsicinterior-eq-interior} gives
    $\operatorname{intrinsicInterior}_{\mathbb R} s' = \operatorname{interior} s'$.
    Substituting yields
    $\operatorname{intrinsicInterior}_{\mathbb R} s
    = \iota_p\,''\,\operatorname{interior} s'$.
\end{proof}

\begin{lemma}[The interior in the direction subspace is nonempty]
    \label{lem:intrinsic-interior-nonempty}
    \lean{LeanEval.ConvexGeometry.interior_translatedSet_nonempty}
    \leanfile{LeanEval/ConvexGeometry/MinkowskiCaratheodory.lean}
    \uses{lem:translated-set-convex, lem:intrinsic-interior-eq-image}
    \leanok
    If $s$ is nonempty, then $\operatorname{interior} s'$ is nonempty.
\end{lemma}
\begin{proof}
    \uses{lem:translated-set-convex, lem:intrinsic-interior-eq-image}
    As $s$ is convex and nonempty, $\operatorname{intrinsicInterior}_{\mathbb R} s$
    is nonempty (\texttt{intrinsicInterior\_nonempty}). By
    Lemma~\ref{lem:intrinsic-interior-eq-image} it equals
    $\iota_p\,''\,\operatorname{interior} s'$, so $\operatorname{interior} s'$ is
    nonempty as well.
\end{proof}

\begin{lemma}[Separation in the direction subspace]
    \label{lem:separation-in-direction}
    \lean{LeanEval.ConvexGeometry.separation_in_direction}
    \leanfile{LeanEval/ConvexGeometry/MinkowskiCaratheodory.lean}
    \uses{lem:translated-set-convex, lem:intrinsic-interior-eq-image}
    \leanok
    Let $y \in s$ with $y \notin \operatorname{intrinsicInterior}_{\mathbb R} s$,
    and put $y' = y - p \in W$. Then $y' \notin \operatorname{interior} s'$, and
    there is a continuous linear functional $g : W \to_{L} \mathbb R$ with
    \[
        g\,w \;<\; g\,y' \qquad \text{for all } w \in \operatorname{interior} s' .
    \]
\end{lemma}
\begin{proof}
    \uses{lem:translated-set-convex, lem:intrinsic-interior-eq-image}
    By Lemma~\ref{lem:intrinsic-interior-eq-image},
    $\operatorname{intrinsicInterior}_{\mathbb R} s
      = \iota_p\,''\,\operatorname{interior} s'$, so
    $y \notin \operatorname{intrinsicInterior}_{\mathbb R} s$ forces
    $y' = y - p \notin \operatorname{interior} s'$. The set
    $\operatorname{interior} s'$ is open (\texttt{isOpen\_interior}) and convex
    (\texttt{Convex.interior}, using Lemma~\ref{lem:translated-set-convex}). Since
    $W$ is a normed space, the geometric Hahn--Banach theorem for an open convex
    set and a point (\texttt{geometric\_hahn\_banach\_open\_point}) yields a
    continuous linear functional $g : W \to_L \mathbb R$ with $g\,w < g\,y'$ for
    every $w \in \operatorname{interior} s'$.
\end{proof}

\begin{lemma}[Extension of the separating functional to $E$]
    \label{lem:extension-to-E}
    \lean{LeanEval.ConvexGeometry.extension_to_E}
    \leanfile{LeanEval/ConvexGeometry/MinkowskiCaratheodory.lean}
    \uses{lem:separation-in-direction}
    \leanok
    Let $g : W \to_L \mathbb R$ be the functional of
    Lemma~\ref{lem:separation-in-direction}. Then there is a continuous linear
    functional $\ell : E \to_L \mathbb R$ whose restriction to $W$ equals $g$, and
    consequently every $z \in s$ with $z - p \in \operatorname{interior} s'$
    satisfies $\ell z < \ell y$.
\end{lemma}
\begin{proof}
    \uses{lem:separation-in-direction}
    By the Hahn--Banach extension theorem
    (\texttt{exists\_extension\_norm\_eq}) applied to $g$ along the isometric
    inclusion $W \hookrightarrow E$, there is $\ell : E \to_L \mathbb R$ with
    $\ell|_W = g$. For $z \in s$ we have $z - p \in W$ and
    $\ell(z - p) = g(z - p)$, and likewise $\ell(y - p) = g\,y'$. Hence
    $z - p \in \operatorname{interior} s'$ gives $g(z - p) < g\,y'$, i.e.\
    $\ell(z - p) < \ell(y - p)$, and adding $\ell p$ to both sides yields
    $\ell z < \ell y$.
\end{proof}

\begin{lemma}[The bound passes to all of $s$]
    \label{lem:bound-passes-to-closure}
    \lean{LeanEval.ConvexGeometry.bound_passes_to_closure}
    \leanfile{LeanEval/ConvexGeometry/MinkowskiCaratheodory.lean}
    \uses{lem:translated-set-convex, lem:intrinsic-interior-nonempty, lem:extension-to-E}
    \leanok
    With $\ell$ as in Lemma~\ref{lem:extension-to-E}, $\ell z \le \ell y$ for all
    $z \in s$.
\end{lemma}
\begin{proof}
    \uses{lem:translated-set-convex, lem:intrinsic-interior-nonempty,
          lem:extension-to-E}
    By Lemma~\ref{lem:extension-to-E} the strict bound $\ell z < \ell y$ holds
    whenever $z - p \in \operatorname{interior} s'$. Since $s'$ is convex
    (Lemma~\ref{lem:translated-set-convex}) with nonempty interior
    (Lemma~\ref{lem:intrinsic-interior-nonempty}),
    \[
        \operatorname{closure}(\operatorname{interior} s')
        = \operatorname{closure} s' \;\supseteq\; s'
    \]
    by \texttt{Convex.closure\_interior\_eq\_closure\_of\_nonempty\_interior}. The
    affine function $z \mapsto \ell z - \ell y$ is continuous and $\le 0$ on the
    image of $\operatorname{interior} s'$ under $\iota_p$, so by
    \texttt{Continuous.le\_on\_closure} it remains $\le 0$ on the closure, which
    contains every $z \in s$ (as $z - p \in s'$). Hence $\ell z \le \ell y$ for all
    $z \in s$.
\end{proof}

\begin{lemma}[The separating functional is non-constant on $s$]
    \label{lem:functional-non-constant}
    \lean{LeanEval.ConvexGeometry.functional_non_constant}
    \leanfile{LeanEval/ConvexGeometry/MinkowskiCaratheodory.lean}
    \uses{lem:intrinsic-interior-nonempty, lem:separation-in-direction, lem:extension-to-E}
    \leanok
    The functional $\ell$ produced in Lemma~\ref{lem:extension-to-E} is
    non-constant on $s$: there exist $a, b \in s$ with $\ell a \ne \ell b$.
\end{lemma}
\begin{proof}
    \uses{lem:intrinsic-interior-nonempty, lem:separation-in-direction,
          lem:extension-to-E}
    By Lemma~\ref{lem:intrinsic-interior-nonempty} pick
    $w \in \operatorname{interior} s'$ and set $a = w + p \in s$. By
    Lemmas~\ref{lem:separation-in-direction} and~\ref{lem:extension-to-E} the
    separation is strict, so $\ell a < \ell y$ with $a, y \in s$, exhibiting two
    distinct values of $\ell$ on $s$.
\end{proof}

\begin{lemma}[Supporting functional at a relative boundary point]
    \label{lem:supporting-functional}
    \lean{LeanEval.ConvexGeometry.supporting_functional}
    \leanfile{LeanEval/ConvexGeometry/MinkowskiCaratheodory.lean}
    \uses{def:exposed-face, lem:separation-in-direction, lem:extension-to-E, lem:bound-passes-to-closure, lem:functional-non-constant}
    \leanok
    Assume $s$ is nonempty. Let $y \in s$ with $y$ \emph{not} in the intrinsic
    interior of $s$. Then there is a continuous linear functional $\ell$,
    non-constant on $s$, such that $y \in F_\ell(s)$, i.e.\
    $\ell z \le \ell y$ for all $z \in s$.
\end{lemma}
\begin{proof}
    \uses{def:exposed-face, lem:separation-in-direction, lem:extension-to-E,
          lem:bound-passes-to-closure, lem:functional-non-constant}
    Work in the direction subspace $W$ with base point $p$ and translated set
    $s'$. Lemma~\ref{lem:separation-in-direction} separates $y' = y - p$ from
    $\operatorname{interior} s'$ by some $g : W \to_L \mathbb R$;
    Lemma~\ref{lem:extension-to-E} extends $g$ to $\ell : E \to_L \mathbb R$, and
    Lemma~\ref{lem:bound-passes-to-closure} shows $\ell z \le \ell y$ for all
    $z \in s$, i.e.\ $y \in F_\ell(s)$ by Definition~\ref{def:exposed-face}.
    Finally Lemma~\ref{lem:functional-non-constant} shows $\ell$ is non-constant on
    $s$.
\end{proof}

\begin{lemma}[A nonzero direction exists]
    \label{lem:exists-direction-vector}
    \lean{LeanEval.ConvexGeometry.exists_direction_vector}
    \leanfile{LeanEval/ConvexGeometry/MinkowskiCaratheodory.lean}
    \leanok
    If $\operatorname{finrank}_{\mathbb R}(\operatorname{vectorSpan}_{\mathbb R} s) \ge 1$,
    then there is a nonzero vector $v \in \operatorname{vectorSpan}_{\mathbb R} s$.
\end{lemma}
\begin{proof}
    A finite-dimensional space whose \texttt{finrank} is at least $1$ is nontrivial
    (\texttt{finrank\_pos\_iff}), hence contains a nonzero element
    (\texttt{exists\_ne}).
\end{proof}

\begin{lemma}[The line section is compact and convex]
    \label{lem:line-section-compact-convex}
    \lean{LeanEval.ConvexGeometry.lineSection_isCompact_convex}
    \leanfile{LeanEval/ConvexGeometry/MinkowskiCaratheodory.lean}
    \leanok
    Let $x \in s$ and let $v \in \operatorname{vectorSpan}_{\mathbb R} s$ be nonzero.
    Then the section
    \[
        T = \{\, t \in \mathbb R : x + t\, v \in s \,\}
    \]
    is compact and convex.
\end{lemma}
\begin{proof}
    The map $\gamma : t \mapsto x + t\, v$ is continuous and affine, and (as
    $v \ne 0$) is a homeomorphism onto its image. Thus $T = \gamma^{-1}(s)$ is
    convex, being the preimage of the convex set $s$ under an affine map
    (\texttt{Convex.affine\_preimage}), and closed, being the preimage of the
    closed set $s$ (\texttt{IsCompact.isClosed}). It is also bounded, since $s$ is
    bounded and $\gamma$ is a homeomorphism onto its image; hence $T$ is closed and
    bounded in $\mathbb R$, i.e.\ compact.
\end{proof}

\begin{lemma}[The line section is a closed interval]
    \label{lem:line-section-icc}
    \lean{LeanEval.ConvexGeometry.lineSection_eq_Icc}
    \leanfile{LeanEval/ConvexGeometry/MinkowskiCaratheodory.lean}
    \uses{lem:line-section-compact-convex}
    \leanok
    With $x$, $v$, and $T$ as in Lemma~\ref{lem:line-section-compact-convex},
    \[
        T = \operatorname{Icc}\!\big(\inf T,\ \sup T\big).
    \]
\end{lemma}
\begin{proof}
    \uses{lem:line-section-compact-convex}
    $T$ is nonempty ($0 \in T$, since $x \in s$), convex and hence connected, and
    closed and bounded (Lemma~\ref{lem:line-section-compact-convex}), so it is
    bounded above and below. A connected, bounded, closed subset of $\mathbb R$
    equals $\operatorname{Icc}(\inf T, \sup T)$
    (\texttt{eq\_Icc\_csInf\_csSup\_of\_connected\_bdd\_closed}).
\end{proof}

\begin{lemma}[Zero is interior to the section]
    \label{lem:line-section-zero-interior}
    \lean{LeanEval.ConvexGeometry.lineSection_zero_interior}
    \leanfile{LeanEval/ConvexGeometry/MinkowskiCaratheodory.lean}
    \uses{lem:line-section-icc}
    \leanok
    With $x$, $v$, and $T$ as above, if $0 \in \operatorname{interior} T$ then
    $\inf T < 0 < \sup T$.
\end{lemma}
\begin{proof}
    \uses{lem:line-section-icc}
    By Lemma~\ref{lem:line-section-icc}, $T = \operatorname{Icc}(\inf T, \sup T)$,
    whose interior is the open interval $\operatorname{Ioo}(\inf T, \sup T)$
    (\texttt{interior\_Icc}). Hence $0 \in \operatorname{interior} T$ means
    $\inf T < 0 < \sup T$.
\end{proof}

\begin{lemma}[Endpoints of the section are not interior]
    \label{lem:segment-endpoints-not-interior}
    \lean{LeanEval.ConvexGeometry.segment_endpoints_not_interior}
    \leanfile{LeanEval/ConvexGeometry/MinkowskiCaratheodory.lean}
    \uses{lem:line-section-icc}
    \leanok
    Let $x \in s$, let $v \in \operatorname{vectorSpan}_{\mathbb R} s$ be nonzero,
    and let $T = \{\, t : x + t\, v \in s \,\}$ with $a = \inf T$ and $b = \sup T$.
    Then neither $x + a\, v$ nor $x + b\, v$ lies in
    $\operatorname{intrinsicInterior}_{\mathbb R} s$.
\end{lemma}
\begin{proof}
    \uses{lem:line-section-icc, lem:line-section-zero-interior}
    If $x + a\, v$ were in the intrinsic interior of $s$, then $0$ would be an
    interior point of the section of $s$ through $x + a\, v$ along $v$, so by
    Lemma~\ref{lem:line-section-zero-interior} (applied at $x + a\, v$) the section
    through $x$ would extend strictly below $a$: some $a' < a$ still has
    $x + a'\, v \in s$, contradicting $a = \inf T$. The same argument at
    $b = \sup T$ rules out $x + b\, v$.
\end{proof}

\begin{lemma}[An interior point lies between two boundary points]
    \label{lem:interior-in-segment}
    \lean{LeanEval.ConvexGeometry.interior_in_segment}
    \leanfile{LeanEval/ConvexGeometry/MinkowskiCaratheodory.lean}
    \uses{lem:exists-direction-vector, lem:line-section-icc, lem:segment-endpoints-not-interior}
    \leanok
    Suppose $\operatorname{finrank}_{\mathbb R}(\operatorname{vectorSpan}_{\mathbb R} s) \ge 1$
    and let $x$ be in the intrinsic interior of $s$. Then there exist $y, z \in s$,
    both \emph{not} in the intrinsic interior of $s$, with $x$ on the segment
    $[y, z]$.
\end{lemma}
\begin{proof}
    \uses{lem:exists-direction-vector, lem:line-section-compact-convex,
          lem:line-section-icc, lem:line-section-zero-interior,
          lem:segment-endpoints-not-interior}
    By Lemma~\ref{lem:exists-direction-vector} pick a nonzero
    $v \in \operatorname{vectorSpan}_{\mathbb R} s$ and form the section
    $T = \{\, t : x + t\, v \in s \,\}$, which equals
    $\operatorname{Icc}(\inf T, \sup T)$
    (Lemmas~\ref{lem:line-section-compact-convex} and~\ref{lem:line-section-icc}).
    Since $x$ lies in the intrinsic interior of $s$ and
    $v \in \operatorname{vectorSpan}_{\mathbb R} s$, a neighbourhood of $0$ stays in
    $T$, so $0 \in \operatorname{interior} T$ and
    Lemma~\ref{lem:line-section-zero-interior} gives $\inf T < 0 < \sup T$. Put
    $a = \inf T$, $b = \sup T$, $y = x + a\, v$, and $z = x + b\, v$, all in $s$.
    Then $x = \tfrac{b}{b-a} y + \tfrac{-a}{b-a} z \in [y, z]$, and by
    Lemma~\ref{lem:segment-endpoints-not-interior} neither $y$ nor $z$ is in the
    intrinsic interior of $s$.
\end{proof}

\subsection{Minkowski's theorem}

\begin{lemma}[Zero-dimensional base case]
    \label{lem:base-case}
    \lean{LeanEval.ConvexGeometry.base_case}
    \leanfile{LeanEval/ConvexGeometry/MinkowskiCaratheodory.lean}
    \leanok
    If $\operatorname{finrank}_{\mathbb R}(\operatorname{vectorSpan}_{\mathbb R} s) = 0$
    and $x \in s$, then $x \in \operatorname{ext}(s)$.
\end{lemma}
\begin{proof}
    A vector span of rank $0$ is trivial, so by
    \texttt{vectorSpan\_eq\_bot\_iff\_subsingleton} the set $s$ is a subsingleton;
    being nonempty (it contains $x$) it equals $\{x\}$. The unique point of a
    singleton is an extreme point.
\end{proof}

\begin{lemma}[Reduction at a relative boundary point]
    \label{lem:minkowski-boundary}
    \lean{LeanEval.ConvexGeometry.minkowski_boundary}
    \leanfile{LeanEval/ConvexGeometry/MinkowskiCaratheodory.lean}
    \uses{def:exposed-face, lem:exposed-face, lem:exposed-face-lower-dim, lem:supporting-functional}
    \leanok
    Assume $s$ is nonempty, and let $y \in s$ not lie in the intrinsic interior of
    $s$. Then there is a set $F \subseteq s$ that is compact and convex, satisfies
    $\operatorname{ext}(F) \subseteq \operatorname{ext}(s)$, contains $y$, and has
    \[
        \operatorname{finrank}_{\mathbb R}\big(\operatorname{vectorSpan}_{\mathbb R} F\big)
        \;<\;
        \operatorname{finrank}_{\mathbb R}\big(\operatorname{vectorSpan}_{\mathbb R} s\big).
    \]
\end{lemma}
\begin{proof}
    \uses{lem:supporting-functional, lem:exposed-face, lem:exposed-face-lower-dim}
    By Lemma~\ref{lem:supporting-functional} there is a functional $\ell$,
    non-constant on $s$, with $y \in F_\ell(s)$. Take $F = F_\ell(s)$. By
    Lemma~\ref{lem:exposed-face} it is compact, convex, contains $y$, and satisfies
    $\operatorname{ext}(F) \subseteq \operatorname{ext}(s)$; by
    Lemma~\ref{lem:exposed-face-lower-dim} its affine dimension is strictly smaller
    than that of $s$.
\end{proof}

\begin{lemma}[Boundary points are convex combinations of extreme points]
    \label{lem:minkowski-boundary-mem}
    \lean{LeanEval.ConvexGeometry.minkowski_boundary_mem}
    \leanfile{LeanEval/ConvexGeometry/MinkowskiCaratheodory.lean}
    \uses{lem:minkowski-boundary}
    \leanok
    Let $s$ be compact convex and assume the induction hypothesis: for every
    compact convex set $s_0$ with
    $\operatorname{finrank}_{\mathbb R}(\operatorname{vectorSpan}_{\mathbb R} s_0)
      < \operatorname{finrank}_{\mathbb R}(\operatorname{vectorSpan}_{\mathbb R} s)$
    and every $x_0 \in s_0$ one has
    $x_0 \in \operatorname{convexHull}_{\mathbb R}(\operatorname{ext}(s_0))$. Then
    for every $y \in s$ not in the intrinsic interior of $s$,
    $y \in \operatorname{convexHull}_{\mathbb R}(\operatorname{ext}(s))$.
\end{lemma}
\begin{proof}
    \uses{lem:minkowski-boundary}
    By Lemma~\ref{lem:minkowski-boundary} there is a compact convex set
    $F \subseteq s$ with $y \in F$,
    $\operatorname{ext}(F) \subseteq \operatorname{ext}(s)$, and
    $\operatorname{finrank}_{\mathbb R}(\operatorname{vectorSpan}_{\mathbb R} F)
      < \operatorname{finrank}_{\mathbb R}(\operatorname{vectorSpan}_{\mathbb R} s)$.
    The induction hypothesis applied to $F$ and $y$ gives
    $y \in \operatorname{convexHull}_{\mathbb R}(\operatorname{ext}(F))
      \subseteq \operatorname{convexHull}_{\mathbb R}(\operatorname{ext}(s))$, the
    inclusion by monotonicity of the convex hull.
\end{proof}

\begin{lemma}[Minkowski's theorem]
    \label{lem:minkowski}
    \lean{LeanEval.ConvexGeometry.minkowski}
    \leanfile{LeanEval/ConvexGeometry/MinkowskiCaratheodory.lean}
    \uses{lem:base-case, lem:minkowski-boundary-mem, lem:interior-in-segment}
    \leanok
    For every compact convex $s$ and every $x \in s$ one has
    $x \in \operatorname{convexHull}_{\mathbb R}(\operatorname{ext}(s))$.
\end{lemma}
\begin{proof}
    \uses{lem:base-case, lem:minkowski-boundary-mem, lem:interior-in-segment}
    We argue by strong induction on
    $d = \operatorname{finrank}_{\mathbb R}(\operatorname{vectorSpan}_{\mathbb R} s)$,
    with the set $s$ (and the point $x$) universally quantified, so that the
    induction hypothesis is exactly the statement for all compact convex sets of
    affine dimension $< d$.

    \emph{Base case $d = 0$.} By Lemma~\ref{lem:base-case},
    $x \in \operatorname{ext}(s)$, and a point lies in the convex hull of any set
    containing it.

    \emph{Inductive step $d \ge 1$.} The induction hypothesis is precisely the
    hypothesis required by Lemma~\ref{lem:minkowski-boundary-mem}, so that lemma
    applies to $s$: every relative boundary point of $s$ lies in
    $\operatorname{convexHull}_{\mathbb R}(\operatorname{ext}(s))$. If $x$ is itself
    a relative boundary point, we are done. Otherwise $x$ lies in the intrinsic
    interior of $s$, and by Lemma~\ref{lem:interior-in-segment} there are relative
    boundary points $y, z \in s$ with $x \in [y, z]$. Applying
    Lemma~\ref{lem:minkowski-boundary-mem} to each of $y$ and $z$ puts both in
    $\operatorname{convexHull}_{\mathbb R}(\operatorname{ext}(s))$; since the convex
    hull is convex, the segment $[y, z]$ is contained in it, and in particular
    $x \in \operatorname{convexHull}_{\mathbb R}(\operatorname{ext}(s))$.
\end{proof}

\subsection{Main theorem}

\begin{theorem}[Minkowski--Carathéodory]
    \label{thm:minkowski-caratheodory}
    \lean{LeanEval.ConvexGeometry.mem_convexHull_finset_extremePoints_of_mem_compact_convex}
    \leanfile{LeanEval/ConvexGeometry/MinkowskiCaratheodory.lean}
    \uses{lem:caratheodory-card, lem:minkowski}
    \leanok
    Let $E$ be a finite-dimensional real normed space, $s \subseteq E$ compact and
    convex, and $x \in s$. Then there is a finite set $t \subseteq \operatorname{ext}(s)$
    with $\#t \le \operatorname{finrank}_{\mathbb R} E + 1$ and
    $x \in \operatorname{convexHull}_{\mathbb R} t$.
\end{theorem}
\begin{proof}
    \uses{lem:caratheodory-card, lem:minkowski}
    By Minkowski's theorem (Lemma~\ref{lem:minkowski}),
    $x \in \operatorname{convexHull}_{\mathbb R}(\operatorname{ext}(s))$. Apply the
    Carathéodory cardinality bound (Lemma~\ref{lem:caratheodory-card}) with
    $A = \operatorname{ext}(s)$: this produces a finite set
    $t \subseteq \operatorname{ext}(s)$ with
    $\#t \le \operatorname{finrank}_{\mathbb R} E + 1$ and
    $x \in \operatorname{convexHull}_{\mathbb R} t$, as required.
\end{proof}

\end{document}
